% Standalone source-record documentation; no research attempts.
\documentclass[11pt]{article}
\usepackage[margin=25mm]{geometry}
\usepackage{fontspec,unicode-math}
\setmainfont{Latin Modern Roman}
\setmathfont{Latin Modern Math}
\usepackage{amsmath,amssymb,mathtools,xurl,hyperref}
\hypersetup{colorlinks=true,urlcolor=blue,linkcolor=blue}
\setcounter{secnumdepth}{0}
\setlength{\parindent}{0pt}
\setlength{\parskip}{5pt}
\setlength{\emergencystretch}{3em}
\newcommand{\sref}[2]{\hyperlink{source:#1:#2}{[S#2]}}
\newcommand{\cataloguescope}[1]{\paragraph{Source record.}#1}
\newcommand{\floor}[1]{\left\lfloor#1\right\rfloor}
\begin{document}
\section{Finite-switch implementation of chattering controls}\label{statement}

\subsection{Definitions and mathematical statement}

\paragraph{Shared notation.}Standard number systems are $\mathbb N,\mathbb Z,\mathbb Q,\mathbb R,\mathbb C$. Specialized notation and index ranges follow the cited source and the definitions below.


A feedback control is derived from state information, often through a value function solving an HJB equation. A chattering control switches infinitely often over a finite time interval. A finite-switch approximation replaces it by finitely many control intervals. Its performance loss compares the resulting cost with the optimal cost for the same system; discretization error additionally depends on the numerical approximation of the value function.
\paragraph{Record scope and qualifications.}
The catalogue title is retained as metadata. Named certificates, contracts, bounds or frameworks in that title are not asserted as results of the original source extract.
The excerpt contains several practical questions but no specified dynamics, switching budget, norm, or rate. Its qualitative finite-switch question is retained without a universal performance bound. The supplied catalogue marks this record partially solved; its present research status was not independently checked.
The source text below is rendered literally where necessary; damaged formulas and composite extracts are not silently repaired into a new mathematical claim.
\paragraph{Statement recorded in the supplied source.}
A.8 Oberman, Adam

I am interested in finite difference methods for high dimensional control and stochastic control problems. The types of general questions I am interested in are in 2 groups: - we know in theory that we can solve the PDE for the value function. How can we do this in practice for medium dimensional (3-4-5) problems? - Given an approximate numerical solution of the PDE for 9

the value function, it isn't trivial to implement a near-optimal control strategy. How do we best do this and can we understand how far from optimal these approximations are? I am curious about how much of a gap there is between theory and practice for nonlinear control problems. Given a control problem, we can write down the PDE for the value function, and we know that this PDE has unique viscosity solutions. From solutions of the PDE, we can then obtain the control. What I would like to know is the following: Suppose we have an approximate solution of the PDE, and we use this approximation to derive a control. How far from optimal will the resulting solution be? Also: suppose the optimal solution swtiches strategy infinitely many times. (This happens even in well-known examples, like stabilizing the linear/nonliner pendulum). Even supposing we had the exact solution, how far off are we from optimal if we approximate the control by one which only switches finitely many times, (which is all that we can do in practice). \sref{statement}{2}
\paragraph{Source status.}
The supplied catalogue labels this source record partially solved; that metadata was checked against the supplied snapshot on 2026-10-11. The label is preserved for provenance and is not an independent verification of present open status or mathematical correctness.
\subsection{Short English statement}
The source asks how approximate value functions and finitely many switches affect control performance.
\subsection{Sources}
\begin{itemize}
\item[\hypertarget{source:statement:1}{S1}] Ulam AI and the UnsolvedMath Contributors, UnsolvedMath: A Curated Collection of Open Mathematics Problems, catalogue entry. Catalogue curation and metadata: CC BY 4.0. The source extract is retained; definitions, English summaries and qualifications are editorial additions.
\url{https://huggingface.co/datasets/ulamai/UnsolvedMath}.
\item[\hypertarget{source:statement:2}{S2}] AIM Workshop Problem Lists, source linked by the catalogue.
\url{https://aimath.org/WWN/optimalcontrol/optimalcontrol.pdf}.
\end{itemize}
\paragraph{Preparation and review.}
The original source and UnsolvedMath contributors retain attribution for the source extract. Codex (GPT-6; exact serving revision not exposed) assisted with definitions, English summaries, scope notes and formatting. No new human mathematical authorship or independently verified solution is asserted. Catalogue identity checked against the complete supplied registry and current repository. Distinct upstream records, including repeated questions and broad or damaged extracts, are retained for complete coverage of the source CS category. A repeated record is not a new independent conjecture. Complete mathematical-equivalence and current-literature review has not been performed. Preparation and final review follow contribution prompts C and D, with source gaps disclosed. Structural and rendering checks do not establish a new theorem or open status.
\label{end:statement}
\end{document}
